% ==================== THEME 2 — AFFECTATIONS / EXPRESSIONS (entraînement) ====================

% ---- aff-egal-vs-double-egal ----
\begin{question}{ent-aff-egal-vs-double-egal}
    Pour affecter la valeur \texttt{18} à la variable \texttt{age}, quelle écriture est correcte ?
    \begin{multicols}{4}
        \begin{reponses}
            \bonne{\checkmark\ \lstinline[language=C]|age = 18;|}
            \mauvaise{\lstinline[language=C]|age == 18;|}
            \mauvaise{\lstinline[language=C]|18 = age;|}
            \mauvaise{\lstinline[language=C]|age := 18;|}
        \end{reponses}
    \end{multicols}
\end{question}

% ---- aff-point-virgule ----
\begin{question}{ent-aff-point-virgule}
    Parmi les lignes suivantes, laquelle est syntaxiquement \textbf{incorrecte} en C ?
    \begin{multicols}{3}
        \begin{reponses}
            \bonne{\checkmark\ \lstinline[language=C]|vitesse = vitesse * 2|}
            \mauvaise{\lstinline[language=C]|vitesse = vitesse * 2;|}
            \mauvaise{\lstinline[language=C]|vitesse = vitesse * 2;;;|}
        \end{reponses}
    \end{multicols}
\end{question}

% ---- aff-incrementation ----
\begin{question}{ent-aff-incrementation}
    Après les lignes suivantes, quelle est la valeur de \texttt{y} ?
    \lstinputlisting[language=c]{sources/codes/ent-incrementation.c}
    \begin{multicols}{4}
        \begin{reponses}
            \mauvaise{6}
            \mauvaise{7}
            \bonne{\checkmark\ 8}
            \mauvaise{ERREUR}
        \end{reponses}
    \end{multicols}
\end{question}

% ---- aff-division-entiere ----
\begin{question}{ent-aff-division-entiere}
    Que vaut \texttt{r} après \lstinline[language=C]|int r = -13 / 4;| ?
    \begin{multicols}{5}
        \begin{reponses}
            \bonne{\checkmark\ -3}
            \mauvaise{-4}
            \mauvaise{-3.25}
            \mauvaise{3}
            \mauvaise{-13}
        \end{reponses}
    \end{multicols}
\end{question}

% ---- aff-modulo ----
\begin{question}{ent-aff-modulo}
    Que vaut \texttt{r} après \texttt{int r = 29 \% 7;} ?
    \begin{multicols}{4}
        \begin{reponses}
            \mauvaise{4.14}
            \mauvaise{4}
            \bonne{\checkmark\ 1}
            \mauvaise{7}
        \end{reponses}
    \end{multicols}
\end{question}

% ---- aff-trace (tableau structuré) ----
\begin{question}{ent-aff-trace}
    Pour chacune des lignes suivantes, indiquer la variable modifiée et sa nouvelle valeur, après exécution.
    \lstinputlisting[language=c]{sources/codes/ent-affectation.c}
    \evaluationProfTableauTrace{2}{%
        \begin{center}
        \begin{tabular}{|c|c|p{4cm}|}
            \hline
            \textbf{Ligne} & \textbf{Variable} & \textbf{Nouvelle valeur} \\
            \hline
            \texttt{num\_i = num\_i + 4;} & \textbf{num\_i} & \textbf{7} \\[0.3cm]
            \hline
            \texttt{num\_f = num\_i / 2;} & \textbf{num\_f} & \textbf{3} \\[0.3cm]
            \hline
            \texttt{num\_f = num\_i / 2.0;} & \textbf{num\_f} & \textbf{3.5} \\[0.3cm]
            \hline
            \texttt{num\_c = 'G';} & \textbf{num\_c} & \textbf{'G'} \\[0.3cm]
            \hline
            \texttt{num\_c++;} & \textbf{num\_c} & \textbf{'H'} \\[0.3cm]
            \hline
            \texttt{num\_i = 21 \% 5;} & \textbf{num\_i} & \textbf{1} \\[0.3cm]
            \hline
            \texttt{num\_i = 24 \% 5;} & \textbf{num\_i} & \textbf{4} \\[0.3cm]
            \hline
        \end{tabular}
        \end{center}
    }
\end{question}
